% =====================================================================
%  Estilo común de los documentos largos del Consejo de Estudiantes
%  Facultad de Ciencias - Universidad de Zaragoza
%
%  Maquetación del «oficio» del Manual de identidad corporativa de la
%  Universidad de Zaragoza (pág. 82), adaptada al Consejo (igual que en
%  actas/acta-consejo.sty): logotipo a 25 mm del borde superior y 10 mm del
%  izquierdo, con 60 mm de ancho; dirección en el margen inferior izquierdo
%  (15 mm); sello de la Facultad de Ciencias en fondo de agua en el margen
%  inferior derecho, recortado según la pág. 60 y tramado al 10 %.
% =====================================================================
\usepackage{iftex}
\ifPDFTeX
  \usepackage[utf8]{inputenc}
  \usepackage[T1]{fontenc}
\else
  \usepackage{fontspec}
\fi
% Tipografía: el Manual usa Myriad Pro, que es de pago y no va en el
% repositorio. Con LuaLaTeX o XeLaTeX, si la carpeta fonts/ tiene
% MyriadPro-Regular.otf, -Bold.otf, -It.otf y -BoldIt.otf, se usa Myriad Pro.
% En otro caso, Source Sans (su pariente libre) o, si tampoco está, Helvetica.
\newif\ifmyriad
\myriadfalse
\ifPDFTeX\else
  \IfFileExists{fonts/MyriadPro-Regular.otf}{\myriadtrue}{}
\fi
\ifmyriad
  \setmainfont{MyriadPro}[
    Path = fonts/, Extension = .otf,
    UprightFont = *-Regular, BoldFont = *-Bold,
    ItalicFont = *-It, BoldItalicFont = *-BoldIt]
  \setsansfont{MyriadPro}[
    Path = fonts/, Extension = .otf,
    UprightFont = *-Regular, BoldFont = *-Bold,
    ItalicFont = *-It, BoldItalicFont = *-BoldIt]
\else
\IfFileExists{sourcesans.sty}{%
  \usepackage[default]{sourcesans}%
}{\IfFileExists{sourcesanspro.sty}{%
  \usepackage[default]{sourcesanspro}%
}{%
  \PackageWarningNoLine{estilo-consejo}{%
    No se encuentra sourcesans.sty: se usa Helvetica}%
  \usepackage[scaled=0.92]{helvet}%
  \renewcommand*{\familydefault}{\sfdefault}%
}}
\fi
\usepackage[spanish,es-noshorthands]{babel}
\usepackage{xcolor}
\usepackage{titlesec}
\usepackage{geometry}
\usepackage{graphicx}
\usepackage{eso-pic}
\usepackage{titling}
\usepackage{fancyhdr}
\usepackage{amssymb}
\usepackage[normalem]{ulem}
\usepackage{enumitem}
\usepackage{tcolorbox}
\tcbuselibrary{breakable}
\usepackage{booktabs}
\usepackage{tabularx}
\usepackage{array}

% --- Colores ---
% DeleAzul: Pantone 653 del Manual (#223D71). DeleRojo: el rojo del logo.
% (Se conservan los nombres para no romper documentos antiguos.)
\definecolor{DeleAzul}{HTML}{223D71}
\definecolor{DeleRojo}{HTML}{EE2722}
\colorlet{azulUZ}{DeleAzul}
\definecolor{DarkGray3}{rgb}{0.3, 0.3, 0.3}
\definecolor{DarkGray4}{rgb}{0.2, 0.2, 0.2}

% --- Fondo: logotipo, dirección y sello en todas las páginas ---
\newcommand{\logodoc}{logo_consejo.pdf}
\newcommand{\sellodoc}{sello_ciencias_agua.pdf}
\newcommand{\direcciondoc}{%
  \textbf{Consejo de Estudiantes}\\
  Facultad de Ciencias\\
  C/ Pedro Cerbuna, 12 · 50009 Zaragoza\\
  cienzdel@unizar.es}
\AddToShipoutPictureBG{%
  \AtPageLowerLeft{%
    \put(\LenToUnit{\paperwidth-155mm},\LenToUnit{-5mm}){%
      \includegraphics[width=155mm]{\sellodoc}}%
    \put(\LenToUnit{15mm},\LenToUnit{15mm}){\makebox(0,0)[bl]{%
      \parbox[b]{60mm}{\raggedright\fontsize{6.5}{8}\selectfont
        \color{azulUZ}\direcciondoc}}}}%
  \AtPageUpperLeft{%
    \put(\LenToUnit{10mm},\LenToUnit{-25mm}){\makebox(0,0)[tl]{%
      \includegraphics[width=60mm]{\logodoc}}}}}

% --- Márgenes ---
% El logotipo acaba a unos 48 mm del borde superior: el texto empieza a 58 mm.
% El bloque de dirección ocupa los 15-28 mm inferiores: el texto acaba a 35 mm.
\geometry{a4paper, top=58mm, bottom=35mm, left=25mm, right=25mm, footskip=14mm}
\setlength{\droptitle}{-1.2cm}

\titleformat{\section}{\color{DeleAzul}\bfseries\fontsize{16pt}{16pt}\selectfont}{\thesection}{1em}{}
\titlespacing{\section}{0pt}{12pt}{8pt}
\titleformat{\subsection}{\color{DarkGray4}\bfseries\fontsize{13pt}{13pt}\selectfont}{\thesubsection}{1em}{}
\titlespacing{\subsection}{0pt}{8pt}{6pt}
\titleformat{\subsubsection}{\bfseries\fontsize{12pt}{12pt}\selectfont}{\thesubsubsection}{1em}{}
\titlespacing{\subsubsection}{0pt}{6pt}{4pt}

\renewcommand{\baselinestretch}{1.0}
\setlength{\parindent}{0pt}
\setlength{\parskip}{6pt}

% --- Pie de página: a la derecha, porque la izquierda es de la dirección ---
\pagestyle{fancy}
\fancyhf{}
\renewcommand{\headrulewidth}{0pt}
\renewcommand{\footrulewidth}{0pt}
\fancyfoot[R]{\footnotesize\color{DarkGray3}\piedepagina\quad\textcolor{DeleRojo}{|}\quad\thepage}
\fancypagestyle{plain}{%
  \fancyhf{}%
  \fancyfoot[R]{\footnotesize\color{DarkGray3}\piedepagina\quad\textcolor{DeleRojo}{|}\quad\thepage}}

% --- Utilidades comunes ---
\newcommand{\va}{$\rightarrow$}
\newcommand{\cuadro}[1]{\textcolor{#1}{\rule[-0.15em]{0.85em}{0.85em}}}
\setlist{itemsep=2pt, topsep=3pt}
